Esta tesis presenta el diseño y la implementación de un prototipo de inteligencia de negocio para integrar procesamiento de datos crediticios, predicción multi-horizonte, gestión de experimentos y visualización analítica en una institución de la economía social y solidaria del Ecuador. El sistema utiliza PostgreSQL para el almacenamiento y el datamart, Apache Airflow para orquestar los flujos de entrenamiento e inferencia, MLflow para registrar experimentos y Apache Superset para los dashboards. La evaluación compara LightGBM, una red neuronal convolucional y un perceptrón multicapa para dieciocho horizontes temporales. En los resultados reportados, LightGBM alcanzó un AUC-ROC promedio de 0,80, mientras que las implementaciones de CNN y MLP mostraron un AUC-ROC de 0,50. El rendimiento de LightGBM disminuyó con el horizonte y el recall descendió hasta 0,034 a dieciocho meses, lo que limita la interpretación operativa de las predicciones de largo plazo. Bajo la configuración evaluada, LightGBM presentó la mayor capacidad discriminativa; sin embargo, la ausencia de repeticiones, pruebas estadísticas, validación externa y una documentación completa de la etiqueta heurística impide afirmar superioridad general o causalidad. La contribución principal es la integración técnica de un flujo de datos, modelado, registro, almacenamiento analítico y visualización, presentada como un prototipo funcional que requiere validación metodológica y operativa adicional.\\

\noindent\textbf{Palabras clave:}

Riesgo crediticio; predicción multi-horizonte; LightGBM; redes neuronales; inteligencia de negocio; MLOps; datamart; Apache Airflow.